Ниже серии расположены по содержательным вопросам, а не по их числовым
идентификаторам: сначала проверяется достижимость границы и роль фазового
порядка, затем параллелизм и декомпозиция, после чего --- масштабная
устойчивость, перенос на синтетические DAG и контролируемое размещение
human-work. Числовой идентификатор обозначает замороженный протокол, а не место
серии в логике изложения.

\subsection{Attainability of \texorpdfstring{$B_4$}{B4}}

Равенство, полученное в EXP-00 для иллюстративного сценария, не является общим
свойством модели. В EXP-01 все 90 сценариев
решены до оптимальности, и в 30 из них найден строгий разрыв $T^*>B_4$;
максимальный относительный разрыв на этой сетке равен 18.333\%. Выбранный
четырёхзадачный контрпример на сетке EXP-01 --- diamond при $P=2$: $B_4=48.00$, но
$T^*=49.25$ часа. Дополнительные 1.25 часа создаёт ресурсная цепочка повторных
человеческих фаз, отсутствующая в трёх агрегатных ветвях $B_4$.

EXP-03 прямо проверяет достаточность агрегатов. В 14 из 24 пар одинаковые DAG,
$A$, $H$, $W_4$, $L_4$, $B_4$, $C$ и $\gamma$ дали разные $T^*$. В наиболее
выраженной выбранной паре при общем $B_4=36.0$ два фазовых профиля дали 40.5 и
37.75 часа. Следовательно, $B_4$ корректно задаёт структурную границу, но
точное значение определяет ресурсно-дополненный фазовый граф.

Рисунок~\ref{fig:resource-augmented-path} показывает такую цепочку явно на
одной точке EXP-05 с $P=2$, $m=4$ и агентной долей overhead $0.05$. В ней
обычная последовательность фаз и зависимости DAG перемежаются порядком
обслуживания единственного человеческого ресурса и порядком задач на агентном
слоте. Поэтому её длина $36.45$ часа совпадает с доказанным $T^*$ и превышает
агрегатную границу $B_4=31.80$ часа.

\begin{figure}[H]
  \centering
  \includegraphics[width=0.96\textwidth]{candidate_resource_augmented_path.png}
  \caption{Ресурсно-дополненный критический путь: task sequence --- порядок
  фаз одной задачи, precedence --- ребро DAG, human order --- порядок на
  человеческом ресурсе, slot order --- порядок задач на одном агентном слоте.
  Обозначения $d1,d2$ относятся к двум декомпозированным подзадачам.}
  \label{fig:resource-augmented-path}
\end{figure}

\subsection{Configured Parallelism}

На 70 точных сценариях EXP-02 при постоянных $C=\gamma=1$ рост доступного $P$
не ухудшал оптимум и быстро приводил к плато по критическому пути или
человеческой ветви. Если же с $P$ росли $C(P)$ или $\gamma(P)$, во всех четырёх
соответствующих исследованных кривых внутренний минимум находился при $P=2$.
Это согласуется с Предложениями~\ref{prop:capacity-monotonicity}
и~\ref{prop:finite-optimum}: ухудшение создаёт не лишний свободный поток, а
изменившаяся конфигурация процесса.

\begin{figure}[H]
  \centering
  \includegraphics[width=0.94\textwidth]{exp02_scaling_curve.png}
  \caption{Точный makespan при изменении числа агентных потоков. Постоянные
  издержки дают насыщение, растущие $C(P)$ или $\gamma(P)$ --- внутренний
  минимум на исследованной сетке.}
  \label{fig:exp02-scaling}
\end{figure}

\subsection{Decomposition and Its Interaction with \texorpdfstring{$P$}{P}}

В EXP-04 точный профиль $T^*(m)$ оказался U-образным в 12 из 13 кривых --- во
всех вариантах с положительными издержками дробления. Общий нулевой срез после
снятия сериализации вышел на плато. При небольшом overhead оптимум обычно
достигался при $m=3$, а при его росте смещался к $m=2$. Поэтому декомпозиция
не монотонно полезна: она сокращает структурную сериализацию, но добавляет
спецификацию, приёмку и интеграцию.

\begin{figure}[H]
  \centering
  \includegraphics[width=0.94\textwidth]{exp04_decomposition_curves.png}
  \caption{Точные кривые декомпозиции. Положительные издержки создают конечную
  оптимальную гранулярность; нулевые издержки дают плато.}
  \label{fig:exp04-decomposition}
\end{figure}

EXP-05 проверяет совместный эффект. Из 100 внутренних interaction contrasts при
положительном overhead 96 положительны, четыре равны нулю и ни один не
отрицателен. Однако абсолютный выигрыш от дробления при $P>1$ положителен лишь
в 74 точках и отрицателен в 26; в 22 точках положительный interaction contrast
сочетается с отрицательным абсолютным выигрышем. Следовательно,
взаимодополняемость означает усиление эффекта относительно $P=1$, а не
безусловную полезность дробления. Увеличение человеческой составляющей overhead
ослабляет выигрыш. Это результат выбранной сетки, а не универсальная теорема.

\subsection{Scale and Parameter Robustness: EXP-06}

EXP-06 воспроизвёл масштабную инвариантность как реализационную проверку на 20
из 20 точных точек.
Преимущество иллюстративного варианта 4 над вариантом 3 сохранилось во всех 5000
случайных и 382 заранее заданных неблагоприятных парах. Этот результат означает
устойчивость лишь внутри замороженных семейств ошибок, а не при произвольном
возмущении всех параметров.

\subsection{Transfer across Synthetic DAGs: EXP-07}

На 1440 сценариях, построенных из 480 синтетических DAG с тремя размещениями
human-work, EXP-07 в пилотном режиме поддержал два заранее заданных медианных
направления для \texttt{bottom\_level\_fcfs\_v1}. В C2 медианный сдвиг
$P_{\mathrm{opt}}:=\min\operatorname*{arg\,min}_{P}T_\pi(P)$ при введении
agent- и human-cost был равен нулю; сдвиги в
противоположную, положительную сторону встретились соответственно в 0.42\% и
0.63\% DAG, то есть ниже предзаданного порога 5\%. В C3 медиана
$I_\pi(4,3)$ равнялась 6.4 с 95\%-м bootstrap-интервалом $[6.0,6.8]$, а
медиана ослабления выигрыша из-за overhead --- 2.0 с интервалом $[2.0,2.0]$.
Отрицательные индивидуальные contrasts при этом составляли 6.53\% и 3.47\%.
Следовательно, пилот поддержал агрегированные направления, а не
универсальные знаки для каждого DAG и не утверждения о доказанном $T^*$.
Из-за частичной реализации pilot-extension rule эти решения относятся к
полученному 20-seed pilot и требуют повторной проверки с полным gate.

Гипотезы о «синхронном» профиле и о простой концентрации human-work не
подтвердились; первая фактически не моделировала синхронизацию во времени, а
вторая смешивала изменение общего $H$ с его расположением. Поэтому
ретроспективная интерпретация первой ограничена сравнением равного и
переставленного асинхронных фазовых профилей, а вывод о расположении $H$
проверяется заново в EXP-08. В exact-аудите EXP-07 только 17 из 36 запусков
получили \texttt{OPTIMAL}; четыре полные пары C1 дали малую оценку
противоположного знака. Этот аудит не охватывал C2--C3 и не заменяет основной
policy-estimand.

\subsection{Placement of Fixed Human Work: EXP-08}

Исправленный парный дизайн EXP-08 сохранил DAG, $A$, $H$ и общий объём работы.
При $C=\gamma=1$ перемещение человеческой работы между задачами не изменило
$H,W_4,L_4,B_4$ и активные ветви во всех 2880 проверках. При $\gamma>C$
перенос доли $\Delta H$ на заранее выбранный путь изменяет его длину точно на
$(\gamma-C)\Delta H$; это тождество подтверждено в 480 из 480 пар. Изменение
же срока фиксированной политики зависело от фаз и $P$. Для первичного профиля
\texttt{io} нормированная парная медиана ``high минус low'' при $P=4$ была
$-0.0077$ с 95\%-м bootstrap-интервалом $[-0.0090,-0.0065]$ и по
предзаданной margin 0.01 классифицирована как practically equivalent. При
$P=8$ оценка составила $-0.0117$ с интервалом $[-0.0120,-0.0101]$ и получила
решение negative.

Все 16 сценариев exact-аудита имели статус \texttt{OPTIMAL}, однако это не
делает статистический вывод определённым: восемь доказанных пар дали для
$P=4$ и $P=8$ одинаковую медиану $0.0010$ с интервалом
$[-0.0150,0.0080]$, поэтому оба решения остались inconclusive. Тем самым
важны две разные характеристики: общий объём человеческой работы задаёт
ресурсный потолок; при $\gamma\ne C$ её размещение меняет веса путей и $L_4$, а
при $C=\gamma$ порядок фаз всё ещё может изменить очередь и срок фиксированной
политики при неизменных агрегатах.
